\chapter{Deduplicación y paginación de caminos}
\label{ch:dedup-pagination}

El matcher puede devolver muchos caminos para una hipótesis laxa.
Antes de que el analista los vea, se deduplican (según el modo
pedido), se puntúan, se filtran, se ordenan y se paginan.
\code{score\_and\_paginate\_paths} es la única función que hace
todo eso; vive en \filepath{graph\_hunter\_core/src/analytics.rs}.

\section{Modos de dedup}

\begin{definition}[Modos de dedup]
\code{DedupMode} es uno de:
\begin{itemize}
    \item \code{None} --- se mantiene cada camino del matcher.
    \item \code{ByPath} --- los caminos con secuencias de
        vértices idénticas se fusionan en un único
        representante (el de mayor score).
    \item \code{ByEndpoints} --- los caminos que comparten el
        mismo par $(v_0, v_\ell)$ se fusionan, sin importar los
        vértices intermedios.
\end{itemize}
\end{definition}

\code{ByPath} es el default para hunts renderizados en la UI:
cuando el grafo subyacente tiene muchas aristas paralelas entre
el mismo par de vértices, el matcher produce múltiples caminos
idénticos a la vista (uno por arista) y el analista sólo quiere
ver la topología una vez.

\code{ByEndpoints} se usa durante análisis en bulk: ``cuántos
pares distintos start-end de alcanzabilidad satisfacen esta
hipótesis'' --- la granularidad correcta para contar escenarios
de ataque únicos.

\section{Algoritmo}

\begin{algorithm}
\caption{\textsc{Score-And-Paginate}}\label{alg:score-paginate}
\KwIn{caminos $\Pi$, índice de página $p$, tamaño de página $s$,
      $\text{min\_score}$ opcional, modo $M$, scorer $S$}
\KwOut{página de caminos + conteo total (pre-página)}
\tcp{Paso 1: dedup según el modo}
\Switch{$M$}{
    \Case{\code{None}}{únicos $\gets \Pi$}
    \Case{\code{ByPath}}{
        clavar cada $\pi \in \Pi$ por su tupla de vértices;
        quedarse con el $\pi$ de mayor score por clave
    }
    \Case{\code{ByEndpoints}}{
        clavar cada $\pi \in \Pi$ por $(v_0, v_\ell)$;
        quedarse con el $\pi$ de mayor score por clave
    }
}
\tcp{Paso 2: puntuar}
\ForEach{$\pi \in $ únicos}{
    $(\text{composite}, \text{breakdown}) \gets S.\text{score\_path}(\pi)$
}
\tcp{Paso 3: filtro opcional}
\If{$\text{min\_score}$ es Some}{descartar $\pi$ con composite $<$ min\_score}
\tcp{Paso 4: sort DESC por (composite, max\_node\_score, total\_score)}
ordenar únicos\;
\tcp{Paso 5: slice}
$\text{total} \gets |\text{únicos}|$\;
\Return (únicos[$p \cdot s$ .. $(p+1) \cdot s$], total)\;
\end{algorithm}

\begin{complexity}
Sea $P = |\Pi|$ (conteo crudo) y $U = |\text{únicos}|$ tras
dedup.
\begin{itemize}
    \item Dedup en \code{ByPath}/\code{ByEndpoints} es
        $\Ord{P \cdot \ell}$ (hashing la clave de largo $\ell$);
        \code{None} es $\Ord{P}$.
    \item Scoring es $\Ord{U \cdot \ell}$ usando \textsc{Score-Path}.
    \item Sort es $\Ord{U \log U}$ comparaciones.
    \item Slicing es $\Ord{s}$.
\end{itemize}
Total: $\Ord{P \cdot \ell + U \log U + s}$. Espacio $\Ord{U + s}$
por encima del input.
\end{complexity}

\section{Estabilidad del sort}

La clave de sort es una tupla
\[
    (\mathrm{composite}(\pi), \max_{v \in \pi}\mathrm{composite}(\{v\}),
     \mathrm{total\_score}(\pi))
\]
ordenada en sentido descendente. La clave de tres vías importa
porque los puntajes compuestos empatan seguido en caminos de
longitud corta --- las claves secundaria y terciaria desempatan
determinísticamente. En particular el max-node score promueve
caminos que contienen un único vértice fuertemente anómalo por
sobre caminos uniformemente mediocres.

\section{Contrato de paginación}

\begin{invariant}[Estabilidad de página]
Dos llamadas consecutivas a \code{get\_hunt\_page} con los mismos
parámetros devuelven el mismo orden, suponiendo que el grafo
subyacente no fue mutado. El orden de caminos es determinista
porque todo desempate es una cantidad totalmente ordenada
derivada del estado del grafo.
\end{invariant}

El arnés de paridad (Capítulo~\ref{ch:parity-harness}) fija este
invariante con una fixture que corre un hunt dos veces y
asevera igualdad elemento a elemento de la página retornada.

\begin{inthecode}
\filepath{graph\_hunter\_core/src/analytics.rs:1096} ---
\code{score\_and\_paginate\_paths}.\\
\filepath{graph\_hunter\_core/src/analytics.rs} --- variantes de
\code{DedupMode} y helpers de derivación de clave.
\end{inthecode}
